\documentclass[10pt,a4paper]{amsart}
\usepackage[margin=19mm]{geometry}
\usepackage{amsmath,amssymb,mathtools}
\usepackage[scr=rsfs,cal=euler]{mathalfa}
\usepackage{xfrac}
\usepackage[unicode,hidelinks,bookmarks=false]{hyperref}
\newcommand{\ie}{i.e.\ }
\newcommand{\cf}{cf.\ }
\newcommand{\resp}{respectively\ }
\newcommand{\Subsection}{\subsection}
\providecommand{\nobreakdash}{\nobreak\hbox{-}}
\setlength{\emergencystretch}{2em}
\multlinegap0pt
\usepackage{bm}
\usepackage{bbm}
\usepackage{dsfont}
\usepackage{xspace}
\usepackage{cancel}
\usepackage{mathtools}
\usepackage{amsthm}
\makeatletter
\@ifundefined{theorem}{\newtheorem{theorem}{Theorem}}{}
\@ifundefined{remark}{\newtheorem{remark}[theorem]{Remark}}{}
\makeatother
\makeatletter
\expandafter\ifx\csname proof\endcsname\relax
\newenvironment{proof}[1][Proof]{\noindent\textbf{#1.} }{\
\rule{0.5em}{0.5em}}
\fi
\makeatother
\title[Necessary Aand Sufficient Characterization Of Absolutely Con]{Necessary Aand Sufficient Characterization Of Absolutely Continuous Functions Defined Over Unbounded Intervals}
\author{Gourav Banerjee}
\date{Source: arXiv:2510.20843v2 (CC BY 4.0)}
\begin{document}
\maketitle

\noindent\emph{Source and licence.} Necessary Aand Sufficient Characterization Of Absolutely Continuous Functions Defined Over Unbounded Intervals. Version arXiv:2510.20843v2 (CC BY 4.0), distributed under the Creative Commons Attribution 4.0 International licence, as recorded in the arXiv licence metadata for this exact version and shown on the abstract page https://arxiv.org/abs/2510.20843v2. Full rights notice: licenses/arxiv-2510.20843-CC-BY-4.0.txt.\par\smallskip

\noindent\textit{Editorial scope.} Counted unit: the Remark whose source label is \texttt{None} in the article (source0.tex lines 360--362, source label \texttt{None}), delivered together with the article's own complete original proof of it (source0.tex lines 364--542, one \texttt{proof} environment of 179 lines). No mathematics is written, repaired or re-derived: statement and proof are the article's own text line for line. Required context retained uncounted: none. Every definition, display and label of those lines is retained unchanged. Omitted from this package and not redistributed: 1--359, 363--363, 543--633. Nothing inside a retained excerpt is omitted or altered: every retained body is delivered exactly as the article's lines are written, so no omission receipt of any kind accompanies this package. Citations: the retained excerpts carry no citation command at all, so no bibliography entry of the article is transcribed and no reference is redistributed. Reference adaptation: none was needed and none was made; every reference inside the delivered unit resolves inside it, so no unresolved reference or citation is produced. Printed slips kept literal: no printed word, symbol, hypothesis or number of the counted statement or of its proof was altered. Layout and class: the article's own class is \texttt{article} with packages including amsmath, amssymb, pgfplots, theorem, geometry, parskip; the wrapper is the lane's pinned \texttt{amsart} editorial wrapper at 10pt on A4, which restates the theorem-like environments the retained text needs and adds the article's own macro definitions that the retained text needs (none), with extra wrapper packages (bm, bbm, dsfont, xspace, cancel, mathtools). The delivered numbering is the unit's own. Assets: no retained span contains a figure environment, a graphics-inclusion command, a PDF-page inclusion, a table or a TikZ picture; no image file is copied into this package, the package declares no asset dependency and no image file is redistributed with it. Licence: version arXiv:2510.20843v2 of this article is distributed under the Creative Commons Attribution 4.0 International licence, as recorded in the arXiv licence metadata for this exact version and shown on the abstract page https://arxiv.org/abs/2510.20843v2. A full rights notice is delivered at licenses/arxiv-2510.20843-CC-BY-4.0.txt. Characters: every retained excerpt is pure ASCII, so the delivered engine, XeLaTeX with its default Latin Modern text font, reports no missing character.
\begin{remark}
A locally absolutely continuous function $f$ is almost everywhere differentiable in $\mathbf{R}$, hence we can talk about its derivative $f'$.
\end{remark}

 \begin{proof}
Let us assume that $f\notin AC(\mathbb{R})$,but$f\in AC_{loc}(\mathbb{R}) $ and  $f' \in L^1_G(\mathbb{R})$, 

Then, $\exists \, \epsilon > 0$, such that for any $\delta > 0$, there exists a sequence of disjoint intervals $(a_i, b_i), \, i \in I$ such that:
\begin{equation}\label{r1}
\sum_{i=1}^{\infty} (b_i - a_i) < \delta \quad \text{but} \quad \sum_{i=1}^{\infty} \left| f(b_i)-f(a_i)  \right| \geq \epsilon.
\end{equation}

Consider a positive sequence $(r_n)$ such that:
\[
\sum_{n=1}^{\infty} r_n < \infty. \quad \text{(In particular, let } r_n = \frac{1}{n^2} \text{.)}
\]

Now, there exists a sequence of intervals $(a_i, b_i)$ such that:
\[
\sum_{i=1}^{k} \left| f(b_i)-f(a_i) \right| > \epsilon \quad \text{but} \quad \sum_{i=1}^{k} (b_i - a_i) < 1.
\]

take $S_1 = \bigcup_{i=1}^{k_1} (a_i, b_i)$, 

Now there exist a closed bounded interval $L_1$
such that $S_1\subset L_1$

Note that $\int_{L_1} |f'| < \infty$ since $f' \in L^1_G$

hence by continuity of measure , given 
\begin{equation}\label{r2}
\epsilon_1 =\frac{\epsilon}{2^2}\,,\,
\exists\,\delta_1 \text{ such that }\int_{E\cap L_1} |f'| < \epsilon_1 
\text{ whenever }\mu(E\cap L_1 ) <\delta_1
\end{equation}

let $m_1 = \min(\delta_1, r_1)$.

Then from (\ref{r1}), there exists $S_2$ (a finite collection of disjoint intervals)
\[
S_2 = \bigcup_{i=1}^{k_2} (a_i, b_i),
\]
such that 
\[
m(S_2) < m_1, \quad \text{and} \quad \sum_{i=1}^{k_2} \left| f(b_{i_2}) - f(a_{i_2})\right| > \epsilon.
\]
since $ f\in AC_{loc}(\mathbb{R}) \implies\int_{a_i}^{b_i} f' d\mu =  f(b_{i_2}) - f(a_{i_2})$
due to Lebesgue's fundamental theorem of calculus 

\begin{equation}\label{r3}
    \implies \sum_{i=1}^{k_2} \left| \int_{a_i}^{b_i} f' \, d\mu \right| > \epsilon.
\implies  \sum_{i=1}^{k_2}  \int_{a_i}^{b_i} \left|f' \right|\, d\mu  > \epsilon.
\implies \int_{S_2} \left|f'\right| >\epsilon
\end{equation}    
\begin{align*}
\implies&&
\int_{S_2 \setminus L_1} \left|f'\right| \, d\mu + \int_{S_2 \cap L_1} \left|f'\right| \, d\mu & > \epsilon\\
\implies && 
\int_{S_2 \setminus L_1} \left|f'\right| \, d\mu &> \epsilon - \int_{S_2 \cap L_1} \left|f'\right| \, d\mu\\ 
\implies&&& >\epsilon-\frac{\epsilon}{2^2}\quad(\because \mu(S_2\cap L_1 <\delta_1 )  \text{ from }\text{(\ref{r2}))}
\end{align*}

then consider $S_2\setminus L_1 = A_2$ and $S_1 = A_1$

In particular, after having chosen $A_n$ , consider a closed bounded interval $L_n $ such that $ \bigcup _{k = 1}^n A_k \subset L_n$, note that $\mu(L_n) <\infty $ therefore $\int_{L_n} \left|f'\right| \, d\mu < \infty$, hence by continuity of  measure , given $\epsilon_n = \frac{\epsilon}{(n+1)^2}$,there exists a $\delta_n$ such that 

\begin{equation}\label{r4}
    \int_{E\cap L_n} \left|f'\right| \, d\mu < \epsilon_n = \frac{\epsilon}{(n+1)^2} \quad \forall \mu(E\cap L_n) <\delta_n
\end{equation}
set $m_n = \min (\delta_n,r_n) $

now from \ref{r1} ,there exists a set $S_{n+1}$(which is a countable collection of disjoint intervals)
such that 
\[\mu(S_{n+1})  \leq m_n \text{  but  } \int_{S_{n+1}} \left|f'\right|\geq\epsilon
\]
\[
\implies \int_{S_{n+1}\setminus L_n} \left|f'\right|\geq\epsilon - \int_{S_{n+1}\cap L_n} \left|f'\right|\geq\epsilon - \frac{\epsilon}{(n+1)^2} \text{    from (\ref{r4})}
\]

hence choose $A_{n+1} = S_{n+1} \setminus L_n$.
Now  we have a countable disjoint collection of open sets $A_n \,,\, n\in \mathbb{N} $ such that 

\[
\mu\left(\bigcup_{n\in \mathbb{N}} A_n \right)\leq
\sum_{n\in \mathbb{N}}\mu(A_n) \leq
\sum_{n\in \mathbb{N}}r_n <\infty
\]
but 
\begin{align*}\displaystyle
    \int_{\bigcup_{n\in \mathbb{N}} A_n } \left|f'\right|d\mu&= \sum_{n = 1}^\infty \int_{ A_n } \left|f'\right|d\mu\\
    &\geq \sum_{n = 1}^\infty \epsilon- \frac{\epsilon}{n^2}\\
    &=\lim_{N\to \infty}N\epsilon - \epsilon\left(\sum_{i = 1}^N \frac{1}{i^2}\right) = \infty
\end{align*}
thereby contradicting the fact that $f' \in L^1_G(\mathbb{R})$\\


\textbf{Converse Part:}

We prove that if $f \in AC(\mathbb{R})$, then $f' \in L^1_G(\mathbb{R})$.  

$Since, f \in AC(\mathbb{R})$, hence for a given $\epsilon = 1$, $\exists \, \delta > 0$, such that if there exists a collection of disjoint intervals $(a_i, b_i)$ such that:
$\sum_{i=1}^\infty (b_i - a_i) < \delta$
then$\sum_{i=1}^\infty \left| f(a_i)-f(b_i)\right| < 1$

Now given any finite measure set $S$, there exists a set of disjoint open intervals (countable) $I_n, \, n \in \mathbb{N}$ such that
\[
S \subset \bigcup_{n=1}^\infty I_n=I, \quad \text{where} \quad \mu\left(\bigcup_{n=1}^\infty I_n\right) = \sum_{n=1}^\infty \mu(I_n) < \infty
\]
Now it suffices to show that:
\[
\int_I |f'| < \infty \text{  which will further imply}  \int_S|f'| < \infty.
\]

now since $\sum_{n=1}^\infty \mu(I_n) < \infty$,
we can have two cases:

\textbf{Case 1: $\sum_{n \in \mathbb{N}} \mu(I_n) < \delta$}

In this case choosing union over finite number of intervals $\bigcup_{k=1}^N I_k$,we can consider any arbitrary partition of this union of intervals ${P_k:k\in\{1,...,N\}}$ such that 
\[
\sum_{k=1}^N t_{a_k}^{b_k}(P_k)<1 \quad \because \sum_{k=1}^n \mu(I_k) < \delta \quad\quad I_k = (a_k,b_k)
\]
here $P_k$ represents an arbitrary partition of the $I_k$$^{th}$ interval. and $t_{a_k}^{b_k}$ represents the variation of $f$ corresponding to $P_k$. Now by taking supremum of variation with respect to all possible $P_k$'s , we have
\[
\sum_{k=1}^N T_{a_k}^{b_k}(f)<1 
\]
here $T_{a_k}^{b_k}(f) = \sup\{t_{a_k}^{b_k}(P_k) \,:\, P_k\text{ is a partition of }(a_k,b_k) \}$ 
we call $\bigcup_{k=1}^\infty I_k = I$
\[
\sum_{k=1}^N T_{a_k}^{b_k}(f)<1 \]
 , but since the function is absolutely continuous over any finite interval
\[
\therefore T_{a_k}^{b_k}(f) =  \int_{I_k} \left|f'\right| \, d\mu 
\]
\[
\implies \sum_{k=1}^N  \int_{I_k} \left|f'\right|  \, d\mu \leq 1,
\]
since this is true for all $N \in \mathbb{N}$
Therefore
\[
\sum_{k=1}^\infty \int_{I_k} |f'|d\mu \implies\int_{I} |f'| d\mu < 1.
\]

\textbf{Case 2: $\sum \mu(I_n) = M > \delta$}

let\[
\bigcup_{n=1}^\infty I_n = I.
\]
now consider $n_0 = \left\lfloor \frac{M}{\frac{\delta}{2}} \right\rfloor$, now note that $M - n_0 \frac{\delta}{2} < \frac{\delta}{2}$.  
also there exist a large finite $N_0 \in \mathbb{N}$, such that:
\[
n_0\frac{\delta}{2} \leq\sum_{n=1}^{N_0} m(I_n) \leq M .
\]

now each of these $N_0$ number of intervals can be subdivided into intervals or union of intervals having length in between $\delta$ and $\frac{\delta}{2},$  let $N'$ be the number of such new disjoint collection of intervals $J_n$, with 
$\frac{\delta}{2}<\mu(J_n)<\delta$ for each $n$.

hence
\[
\int_{J_n} |f'| = \text{Total Variation}(J_n) \leq 1, \quad \forall n\in \mathbb{N}.
\]

Now
\begin{equation}\label{n1}
\int_{\bigcup_{n=1}^{N'} J_n} |f'| d\mu\leq N'\leq n_0
\end{equation}
also $I-\bigcup_{n=1}^{N'} J_n$ has measure $<$ $\frac{\delta}{2}$
hence by proof of case-I
\begin{equation}\label{n2}
\int_{I\setminus\bigcup_{n=1}^{N'} J_n} |f'| d\mu\leq 1
\end{equation}

therefore from (\ref{n1}) and (\ref{n2}) , we get 
\[
\therefore \int_{I} |f'| d\mu\leq n_0+1
\]
\[
\therefore \int_{S} |f'| d\mu\leq n_0+1
\]
hence finite.
\\
Without loss of generality, we can conclude that if  $f\in AC_{loc}(\mathbb{I})$ and $f' \in L^1_G(\mathbb{I}) \iff f\in AC(\mathbb{I}) $ where $\mathbb{I}$ is an interval in $\mathbb{R}$ .
\end{proof}

\end{document}
